\section{Related Work}\label{sec:related}

\subsection{Rough Sets, Reducts, and Attribute Reduction}
 
The closest classical neighbor is rough-set theory and its large literature on reducts and attribute reduction~\cite{pawlak1982rough,jensen2004semantics,jensen2008rough}. That literature asks when a smaller attribute set preserves discernibility or decision-table structure, and it has developed both exact and heuristic methods for reduct computation.

\subsection{Backdoor Tractability in CSPs}
A closely related line of work studies backdoor tractability for constraint satisfaction problems~\cite{williams2003backdoors, bessiere2013detecting}. In that framework, a \emph{backdoor} is a minimal variable set whose fixation collapses a CSP to a tractable class (e.g., one admitting a majority polymorphism). Bessiere et al.~\cite{bessiere2013detecting} showed that testing membership in tractable classes characterized by majority or conservative Malt'sev polymorphisms is polynomial-time, but finding a minimal backdoor to such a class is W[2]-hard when the tractable subset is unknown and fixed-parameter tractable only when that subset is given as input. This directly parallels our structural results: coordinate sufficiency in decision problems and backdoor tractability in CSPs both formalize the same intuition, that discovering what matters is fundamentally harder than exploiting a known structure. Our bounded-treewidth reduction (Theorem~\ref{thm:treewidth}) explicitly connects to their CSP algorithms and backdoor methodology. Subsequent work in parameterized complexity has rigorously bounded the tractability of discovering these minimal structural sets across varying graph widths and backdoor types \cite{ganian2017discovering}.

The complexity of computing minimal reducts has been studied in the rough-set literature. Finding a minimal attribute set that preserves decision-distinctions is known to be NP-hard in general~\cite{slowinski1995decision}, and more refined analysis places certain reduct variants in $\Sigma_2^P$~\cite{multiplier}. These results establish that the static regime questions studied here, whether a coordinate set preserves decision distinctions, and whether a minimal such set exists, fall within a known complexity landscape for attribute reduction.

The key structural difference is this: general decision tables can admit multiple incomparable minimal reducts, necessitating a combinatorial search through the reduct lattice. Under the exact optimal-action preservation predicate studied here, Proposition~\ref{prop:minimal-relevant-equiv} proves there is exactly one minimal sufficient set: the relevant-coordinate set $R(\mathcal{D})$ itself. This uniqueness follows from the optimizer quotient structure: states in the same quotient class share the same optimal-action set, so any sufficient set must separate all quotient classes, and the relevant-coordinate set is precisely the minimal set achieving this separation.

The complexity consequence is significant. In general rough-set reduct computation, multiple minimal reducts can exist, placing minimal-reduct search in $\Sigma_2^P$. Here, uniqueness collapses the search problem: finding a minimal sufficient set reduces to checking which coordinates are relevant (coNP-complete). The $\Sigma_2^P$ search layer vanishes because there is no need to search among incomparable candidates; the single minimal set is uniquely determined by relevance. This structural collapse—from $\Sigma_2^P$ search to coNP relevance containment—is driven by optimizer-quotient uniqueness and is, to our knowledge, a novel classification not captured by general rough-set bounds.

Both settings ask which coordinates can be removed without losing decision distinctions.

\leanmetapending{\LHrng{DP}{1}{2}}
\begin{proposition}[Static Sufficiency as Reduct Preservation for the Induced Decision Table]\label{prop:static-rough-set-bridge}
Given a finite decision problem $\mathcal{D}=(A,X_1,\ldots,X_n,U)$, form the induced decision table whose condition attributes are the coordinates $X_i$ and whose decision label at state $s$ is the set-valued class label $\Opt(s) \subseteq A$. Then a coordinate set $I$ is sufficient in the sense of this paper if and only if the projection to $I$ preserves the decision label of the induced table. Consequently, the minimal sufficient sets of $\mathcal{D}$ are exactly the reducts of this induced decision table.
\end{proposition}


\begin{proof}
By definition, $I$ is sufficient exactly when any two states agreeing on $I$ have the same optimal-action set. In the induced decision table, the decision label of state $s$ is precisely $\Opt(s)$, so the same condition says that projection to $I$ preserves the decision label. Minimality on the present side is therefore exactly reduct minimality for the induced table.
\end{proof}

Rough-set reducts are formulated relative to indiscernibility relations or decision tables, whereas the object here is the optimizer map of a decision problem and the induced exact preservation of optimal-action sets. Proposition~\ref{prop:static-rough-set-bridge} shows that the static regime can be recast as reduct preservation for the induced decision table. What the present paper adds is: (i) the collapse theorem showing minimum sufficiency collapses to relevance containment, bypassing the general NP-hardness barrier for exact minimization; and (ii) the systematic extension to stochastic and sequential regimes, which takes the reduct question beyond the classical rough-set setting. The stochastic and sequential layers introduce conditional-comparison and temporal-structure phenomena that are outside the usual rough-set formulation.

The same generosity is due to the broader feature-selection and subset-selection literature in machine learning and combinatorial optimization~\cite{blum1997selection,amaldi1998complexity}. Those works study sparse predictive representations, informative feature subsets, and the computational cost of choosing them. Our setting differs in targeting exact preservation of the optimal-action correspondence rather than prediction loss, classifier complexity, or approximate empirical utility. Still, the shared concern is clear: identifying what information matters and what can be safely omitted.

\subsection{Decision-Theoretic Sufficiency, Informativeness, and Value of Information}

Classical decision theory and informativeness ask when one information structure is at least as useful as another~\cite{blackwell1953equivalent,savage1954foundations,raiffa1961applied,howard1966information}. Our setting is the computational certification layer of that question: exact coordinate projection preserving optimal-action correspondence.

\subsection{Abstraction, Bisimulation, and Homomorphisms in Planning and Reinforcement Learning}

The stochastic and sequential parts of the paper sit near a second major literature: exact abstraction in Markov decision processes, POMDPs, bisimulation-based aggregation, and MDP homomorphisms~\cite{papadimitriou1987mdp,littman1998probplanning,mundhenk2000mdp,givan2003equivalence,ravindran2004algebraic}. This body of work studies when state spaces can be aggregated while preserving value or policy structure, and it has established many of the representation-sensitive complexity phenomena that make richer stochastic and sequential models computationally difficult.

That literature is structurally close to the present paper and should be credited as such. The main difference is that our predicate is coordinate sufficiency for preserving optimal-action sets under an explicit coordinate-hiding operation. We do not study arbitrary abstract state maps, approximate value-function guarantees, or the quantitative bisimulation metrics pioneered by Ferns et al. \cite{ferns2004metrics} to measure behavioral similarity. While recent work has successfully leveraged such metrics to learn robust, task-invariant representations in high-dimensional deep reinforcement learning (e.g., Zhang et al. \cite{zhang2021learning}), our focus remains strictly on the combinatorial complexity of exact certification. Instead of optimizing a continuous representation space to discard task-irrelevant distractors, we ask whether one can exactly certify that no omitted coordinate changes the optimal-action correspondence. The contribution here is therefore not abstraction theory in general, but a unified exact-certification account in which static, stochastic, and sequential regimes are analyzed as variants of one decision-relevance question.

\subsection{Explanatory and Structural Reasoning Links}

Anchor-style explanations, abductive explanations, and exact reason/justification frameworks all study small structures sufficient for a decision outcome~\cite{ribeiro2018anchors,ignatiev2019abduction,marques2022delivering,darwiche2023reasons}. Our contribution is not a new explanation semantics but a regime-sensitive complexity classification for exact coordinate-preservation predicates.

The same applies to causality/responsibility traditions~\cite{pearl2009causality,spirtes2000causation,halpern2001explanations,chockler2004responsibility}: they motivate relevance notions, while this paper studies the algorithmic certification cost of exact preservation under coordinate hiding.

\subsection{Oracle, Query, and Access-Based Lower Bounds}

Query-access lower bounds provide another nearby technique family~\cite{dobzinski2012query}. Our witness-checking duality and access-obstruction results belong to that tradition in spirit. The difference is that the lower bounds are developed around the same fixed sufficiency predicate used throughout the paper, which allows direct comparison between forward evaluation, certification, exact minimization, and restricted-access models without changing the underlying decision relation.

\subsection{Scope and Novelty}

The paper is intentionally synthetic. Rough sets and reducts study decision-preserving attribute reduction; feature selection studies informative subsets; statistical decision theory studies informativeness and sufficiency; abstraction and bisimulation literatures study exact aggregation of stochastic and sequential decision processes; explanation methods study local decision-preserving cores. The contribution here is to treat exact relevance certification itself as the common object and to develop one coherent complexity-theoretic account around it.

What is distinctive in the present paper is not an isolated static or sequential hardness theorem taken on its own. The static regime can be recast as reduct preservation for the induced decision table, and the sequential row belongs near existing exact abstraction phenomena in richer control models. The contribution here is to formalize one exact decision-preservation predicate in optimizer language, track it coherently across static, stochastic, and sequential regimes, expose the preservation/decisiveness split in the stochastic case, and identify the tractable and intractable boundaries of exact certification. The paper's novelty is therefore best understood as a unified exact-certification theory rather than as a single isolated theorem.
